% 本文件由 LaTeX 排版 Agent 根据有效 translation.json 自动生成。
% author=Tertullian; work_id=0324; title=圣殉道者佩尔佩图阿与费利西塔斯受难记

\chapter{圣殉道者佩尔佩图阿与费利西塔斯受难记}

% structure: p0001 source=h1 target=omitted reason=page-title-already-rendered-by-parent

% structure: p0002 source=h2 target=section reason=relative-heading-rank; base=h2

\section{前言}

若昔日的信德典范——既见证天主的恩宠，又助人获得造就——被收集成文，使人通过阅读它们，如同事实再现，从而尊崇天主，坚强人灵；那么，为何不也收集新的典范，使之同样合于双重目的呢？——即使仅仅因为这些现代的榜样终将成为古旧，并为后世所用，尽管在当代，因着对古代的虚妄尊崇，它们被认为权威较小。但让人自己判断吧，如果他们认为圣神的能力是同一的，不分时期与季节；既然有些较晚的事迹，因临近最后时刻，按照末世所彰显的恩宠丰盈，应被更加重视。因为主说：\Quote{在末日，我要将我的圣神倾注在一切有血肉的人身上；你们的儿女要说预言。我要将我的圣神倾注在我的仆人和使女身上；你们的青年要见异象，你们的老人要做异梦。} 因此，我们——既承认并尊敬现代异象，如同预言一样应许给我们，又视圣神其他的能力为教会的行动，祂为此被派遣，在众人中分施各项神恩，正如主按自己的意愿分给每人——也当必要地将它们收集成文，并藉阅读纪念它们，以光荣天主；这样，任何信德的软弱或沮丧就不会以为天主的恩宠只存在于古人身上，无论是在提拔殉道者的屈尊俯就中，还是在赐予启示中；因为天主始终实现祂所应许的，为不信者作证，为信者谋益。因此，我们将所听见、所摸到的，也传报给你们，弟兄们和小孩子们，使你们这些亲身经历此事的人，能再次纪念以光荣主，也使你们这些仅听闻的人，能与蒙福的殉道者共融，并藉着他们与主耶稣基督共融。愿光荣与尊威归于祂，直到永远。阿们。\par

% structure: p0004 source=h2 target=section reason=relative-heading-rank; base=h2

\section{第一章 论据——当圣人们被逮捕时，圣佩尔佩图阿成功抵抗了父亲的恳求，与其他人一同受洗，被投入肮脏的地牢。她因挂虑幼子，藉着赐予她的异象，得知她的殉道很快就要发生}

1. 年轻的\OriginalTerm{慕道者}{catechumen}雷沃卡图斯及其同伴费利西塔斯、萨图尔尼努斯和塞孔杜努斯被逮捕。其中还有维维亚·佩尔佩图阿，出身高贵，受过良好教育，已出嫁，有一父一母和两个兄弟，其中一个与她一样是慕道者，还有一个尚在哺乳的幼子。她本人约二十二岁。从此刻起，她将亲自叙述自己殉道的全部过程，正如她用自己手笔和心智所留下的记述。\par

2. \Quote{当我们仍与迫害者在一起时，}她说，\Quote{我的父亲，出于对我的爱，坚持试图使我背离，并使我从信德上跌倒。我对他说：\InnerQuote{父亲，你看，比方说，这里放着这个器皿，是个小罐子，还是别的？}他说：\InnerQuote{我看见它是这样。}我回答他说：\InnerQuote{它能被称作别的名字，而不是它本身吗？}他说：\InnerQuote{不能。}\InnerQuote{我也不能称自己为别的什么，只能是我所是的：基督徒。}于是，我的父亲被这话激怒，扑向我，仿佛要挖出我的眼睛。但他只是使我痛苦，然后被魔鬼的诡计击败，离开了。此后，在没有父亲的几天里，我感谢了主；他的缺席成了我的安慰。就在那几天里，我们受了洗，圣神指示我，在\Emph{圣洗}的水中除了祈求身体的坚忍外，别无他求。过了几天，我们被带进地牢，我非常害怕，因为我从未经历过这样的黑暗。可怕的日子！由于人群，兵士冲撞的酷热！我因挂念幼子而异常痛苦。当时有泰尔提乌斯和蓬波尼乌斯，那些蒙福的\OriginalTerm{执事}{deacon}，他们服侍我们，并已通过贿赂安排我们出去到监狱中较舒适的地方休息几小时。于是我们走出地牢，各自照顾自己的需要。我哺乳我的孩子，他因饥饿已虚弱不堪。我因挂虑他，对母亲说话，安慰我的兄弟，并将我的儿子托付给他们照料。我萎靡不振，因为我看见他们因我而憔悴。我忍受这样的忧虑许多天，并获准让我的婴儿与我一同留在地牢里；我立刻强壮起来，对婴儿的忧虑和挂念得以解除；地牢对我来说如同宫殿，我宁愿待在那里，也不愿在其他地方。}\par

那时，我的兄弟对我说：\InnerQuote{亲爱的姐姐，你已处于极高的尊位，你可以求一个神视，好让你知道这次的结果是受难还是得脱。}我自知能蒙恩与主交谈，祂的仁慈我早已体验极深，便勇敢地答应他，说：\InnerQuote{明日我告诉你。}我就祈求，所指示我的就是这些。我看见一架金梯，高不可测，直通到天上，但极窄，只能一人一梯而上；梯子两侧布满了各种铁制武器，有剑、矛、钩、匕首；若有人不小心或未抬头向上，便会被撕碎，血肉粘在铁器上。梯子下面蜷伏着一条巨硕无比的龙，它埋伏着，等待那些登梯的人，恐吓他们不让他们上去。\Person{撒图鲁斯}第一个上去，他后来为我们自愿交出自己，被捕时并不在场。他到达梯顶后，转身对我说：\InnerQuote{\Person{裴蓓}，我在等你；但要小心，免得那龙咬你。}我说：\InnerQuote{因主耶稣基督之名，它不能伤害我。}那龙从梯子下面，仿佛怕我似的，慢慢抬起头来；我踏上第一级时，就踩了它的头。我上去后，看见一片广大的花园，园中有一位白发者，牧人装束，身材高大，正在挤羊奶；周围站着成千上万穿白衣的人。他抬起头，注视着我，对我说：\InnerQuote{欢迎你，女儿。}他叫我，从他挤的奶酪中给了我像一小饼的东西，我合手接过来，吃了它；所有站着的人都说\InnerQuote{阿们}。听到他们的声音，我醒了过来，口中还留有无法形容的甘甜。我立即将这事告诉了我的兄弟，我们明白了这将是受难，从此不再对这个世界存有任何希望。\par

% structure: p0008 source=h2 target=section reason=relative-heading-rank; base=h2

\section{第二章。梗概。\Person{裴蓓}遭父亲围攻，安慰他。被押至法庭时，她自认是基督徒，与其他同伴一同被判喂野兽。她为已故的弟弟\Person{狄诺克拉特}祈祷。}

几天后，传言说我们就要受审。那时我父亲从城里来，焦虑万分。他来到我面前，想要使我跌倒，说：\InnerQuote{可怜我的白发吧，女儿。可怜你的父亲吧，如果我还配称为你的父亲。若我用这双手把你抚养到这如花的年纪，若我曾把你置于你所有的兄弟之上，就不要使我受人的嘲笑。顾念你的兄弟，顾念你的母亲和姨母，顾念你的儿子，他离了你就不能活。放下你的勇气，不要使我们全都毁灭；因为若你遭遇什么，我们中没有一个人能自由说话。}父亲带着深情说这些话，吻我的手，伏在我脚下；他流着泪，不叫我女儿，而叫我女士。我为父亲的白发而悲伤，因为全家只有他不愿因我的受难而喜乐。我安慰他说：\InnerQuote{在那刑架上，天主愿意什么就发生什么。要知道，我们并非在自己手中，而是在天主手中。}于是他悲伤地离开了我。\par

又有一天，我们正在吃饭时，突然被带走受审，来到了市政厅。消息立刻在公共场所附近传开，聚集了无数的人。我们登上台。其余的人被审问，都承认了。然后他们来到我面前，我父亲立即带着我的孩子出现，把我从台阶上拉下来，用恳求的语气说：\InnerQuote{可怜你的孩子吧。}代理总督\Person{希拉里亚努斯}刚刚接替已故的前总督\Person{米努齐乌斯·提米尼亚努斯}获得生杀大权，他说：\InnerQuote{饶恕你父亲的白发，饶恕你孩子的幼年吧，为皇帝们的福泽献祭。}我回答说：\InnerQuote{我不这么做。}\Person{希拉里亚努斯}说：\InnerQuote{你是基督徒吗？}我回答：\InnerQuote{我是基督徒。}我父亲站在那里想使我 \Emph{从信仰上} 跌倒，\Person{希拉里亚努斯}下令把他扔下去，用棍棒打他。父亲的遭遇使我悲伤，如同我自己挨打一样，我为他的可怜暮年而哀痛。于是代理总督对我们所有人宣判，判我们喂野兽，我们高高兴兴地走进地牢。那时，我的孩子习惯吃我的奶，也习惯与我一同在狱中，所以我打发执事\Person{庞波尼乌斯}去我父亲那里要孩子，但父亲不肯给他。然而，正如天主所愿，孩子不再想吃奶，我的乳房也没有让我感到不适，免得我同时因挂虑孩子和乳房的痛苦而受折磨。\par

3. 几天后，当我们都在祈祷时，突然，在祈祷中间，有一句话来到我心中，我提起了\OriginalTerm{迪诺克拉特斯}{Dinocrates}；我惊讶那名字此前从未进入我的脑海，我因想起他的不幸而悲伤。我随即感到自己堪当，并受召为他求情。于是我为他开始恳切祈祷，向主呻吟哀哭。当晚，毫不迟延，在异象中这事启示给我。我看见\Person{迪诺克拉特斯}从一个幽暗之处出来，那里还有另外几个人，他干渴焦灼，面容污秽，肤色苍白，脸上带着他死时的那伤口。这\Person{迪诺克拉特斯}是我的同胞兄弟，七岁时患病惨死——他的脸被癌肿吞噬殆尽，以致他的死亡令众人厌恶。我曾为他祈祷，在我与他之间隔着一段宽大的距离，我们彼此无法靠近。而且，在同一地方，有一池满水，池沿高过男孩的身高；\Person{迪诺克拉特斯}挺起身子，仿佛要喝。我悲伤的是，尽管池中有水，但因池沿太高，他无法喝到。我醒来，知道我的兄弟在受苦。但我相信我的祈祷能减轻他的痛苦；于是我每天为他祈祷，直到我们转入军营监狱，因为我们要在营中竞技场上搏斗。那时是\OriginalTerm{格塔·凯撒}{Geta Cæsar}的诞辰，我日夜为我的兄弟祈祷，呻吟流泪，恳求他得蒙恩准。\par

4. 然后，在我们戴着镣铐的那天，这事显示给我。我看见那地方，先前我见它幽暗，现在却明亮；\Person{迪诺克拉特斯}身体洁净，衣着整齐，正在休息。伤口处我只见伤疤；那池子，我先前所见，如今看见池沿降低到男孩的肚脐，\Emph{如今看见}。有人不停地从池中取水，池沿上有一个盛满水的杯子；\Person{迪诺克拉特斯}走近开始饮用，杯子总不枯竭。当他喝足了，他便离开水边，像儿童一样欢乐地玩耍，我便醒来。于是我明白，他已被从惩罚之地转移。\par

% structure: p0013 source=h2 target=section reason=relative-heading-rank; base=h2

\section{第三章 提要。佩尔佩图阿再次受父亲诱惑。她的第三个异象，在其中她被带去与一个埃及人搏斗。她战斗、得胜，并获得奖赏。}

1. 又过了几天，一个名叫\Person{普登斯}的士兵，监狱的助理监督，开始很尊重我们，看出天主的大能在我们内，便允许许多弟兄来探视我们，好让我们和他们彼此获得慰藉。当竞技的日子临近时，我父亲因痛苦而憔悴，来到我面前，开始拔自己的胡须，扑倒在地，俯伏于地，咒骂自己的年岁，说出那些足以感动天地的话语。我为他不幸的晚年而悲伤。\par

\Quote{我们在要战斗的前一天，我在异象中看见执事\Person{庞波尼乌斯}来到监狱门口，用力敲门。我出去给他开门；他穿着一件装饰精美的白袍，袍上有多重\Emph{calliculae}。他对我说：\InnerQuote{佩尔佩图阿，我们在等你；来！}他向我伸手，我们开始走过崎岖蜿蜒的地方。好容易才气喘吁吁地到达竞技场，他领我进入场中央，对我说：\InnerQuote{不要害怕，我与你同在，并与你一同劳苦；}然后他离开了。我诧异地看着无数的人群。因为我知道我已被交给野兽，但我惊讶野兽没有向我放出。这时一个埃及人出来向我挑战，样貌可怖，有帮手同来，要与我搏斗。又有些俊美的青年前来作为我的帮助和鼓励；我被脱去衣服，变成了男人。然后我的帮手开始用油擦我，如同竞赛的惯例；我看见那埃及人却在另一边滚在尘土里。又有一个人出来，身材奇高，甚至超过竞技场顶端；他穿着一件宽松的束腰外衣，两带之间胸口正中有一件紫色长袍；他穿着各种形状的\Emph{calliculae}，由金银制成；他拿着一根棍子，如同角斗士教练一般，还有一根绿枝，上面有金苹果。他叫大家安静，说：\InnerQuote{这埃及人若胜过这女子，就用剑杀死她；若她战胜他，她就得到这枝条。}然后他离开了。我们彼此靠近，开始互相攻击。他试图抓住我的脚，我用脚后跟踢他的脸；我被提到空中，开始这样踢他，仿佛在践踏大地。但当我看到有些耽搁时，我双手交握，手指相交；我抓住他的头，他脸朝下跌倒，我踩在他头上。人们开始欢呼，我的帮手也欣喜。我走近教练，接过枝条；他亲吻我，对我说：\InnerQuote{女儿，愿平安与你同在。}我开始荣耀地走向\Place{萨纳维瓦里亚门}。于是我醒来，明白我不是要与野兽搏斗，而是与魔鬼搏斗。然而我知道胜利已等着我。这些，迄今为止，我在竞技前几日已完成；但竞技本身发生的事情，谁愿意谁去写吧。}\par

% structure: p0016 source=h2 target=section reason=relative-heading-rank; base=h2

\section{第四章 提要。萨图鲁斯在异象中，与佩尔佩图阿被天使带入大光之中，看见殉道者。被带至天主宝座前，接受亲吻。他们使主教奥普塔图斯和司铎阿斯帕修斯和好。}

1. 此外，蒙福的\Person{撒杜鲁斯}也记述了他所见的这个异象，是他亲自写下来的：——\Quote{我们受苦了}，他说，我们便离开了肉身，开始被四位天使托着向东飞去；他们的手没有触摸我们。我们不是仰面漂浮着向上看，而是像在爬一个缓坡。我们被释放后，终于看到了那最初的无尽光明；我说，\Quote{佩尔佩图阿}（因为她在我的旁边），\Quote{这就是主应许给我们的；我们已领受了应许。}当我们正被那同四位天使托着时，我们面前出现了一个广阔的空间，像一座游乐园，里面有玫瑰树和各种花卉。树的高度如柏树，叶子不断落下。此外，在游乐园里又出现了四位天使，比先前那些更明亮；他们看见我们，便表示尊敬，并对其余的天使赞叹地说：\Quote{他们来了！他们来了！}那托着我们的四位天使大为恐惧，把我们放下；我们步行穿过一条宽路，走了一弗隆长的路程。在那里，我们遇到了\Person{约昆杜斯}和\Person{撒图尔尼努斯}和\Person{阿塔克西乌斯}，他们遭受了同样的迫害，被活活烧死；还有\Person{昆图斯}，他本身也是殉道者，死在监狱里。我们问他们其余的人在哪里。天使们对我们说：\Quote{先来，进去问候你们的主。}\par

2. 我们走近一个地方，其墙壁仿佛是用光建造的；在门口站着四位天使，他们给进去的人穿上白袍。我们穿好之后，进入里面，看到了无尽的光明，并听到一些人齐声不停地呼喊：\Quote{圣！圣！圣！}在中间，我们看到仿佛一位老人坐着，头发雪白，面容年轻；我们看不见他的脚。在他右边和左边有二十四位长老，他们后面还站着许多其他人。我们极其惊奇地进去，站在宝座前；四位天使扶起我们，我们亲吻了祂，祂用手抚摸我们的脸。其余的长老对我们说：\Quote{我们站起来吧；}我们便站起来彼此问候平安。长老们对我们说：\Quote{去享受吧。}我说：\Quote{佩尔佩图阿，你已经得到了你所愿的。}她对我说：\Quote{感谢天主，因为我在肉身中虽已喜乐，如今在这里更加喜乐。}\par

3. \Quote{我们走了出来，看见入口处右边站着主教\Person{奥普塔图斯}，左边站着司铎兼教师\Person{阿斯帕西乌斯}，他们分开放着，面容悲伤；他们俯伏在我们脚前，对我们说：\InnerQuote{请在我们之间恢复和平，因为你们已经出去，留下了我们这样。}我们对他们说：\InnerQuote{你们不是我们的父亲，你是我们的司铎，怎该俯伏在我们脚前呢？}我们也俯伏在地，拥抱他们；\Person{佩尔佩图阿}开始与他们交谈，我们把他们拉到游乐园的玫瑰树下。我们正与他们说话时，天使们对他们说：\InnerQuote{让他们休息吧；如果你们之间有什么不和，要彼此宽恕。}天使们就把他们赶走了。天使们对\Person{奥普塔图斯}说：\InnerQuote{责备你的群众，因为他们聚集到你那里，仿佛从竞技场回来，争论派别问题。}然后我们觉得他们好像要关门了。在那里，我们开始认出许多弟兄，还有许多殉道者。我们都因一种说不出的香气而饱足。然后，我喜乐地醒了。}\par

% structure: p0020 source=h2 target=section reason=relative-heading-rank; base=h2

\section{第五章 论题。塞昆杜卢斯死于狱中。斐利奇达怀孕，但经多次祈祷，在第八个月无痛苦生产，佩尔佩图阿与撒杜鲁斯的勇气未衰。}

1. 以上是蒙福的殉道者\Person{撒杜鲁斯}和\Person{佩尔佩图阿}亲自记述的较为卓越的异象。但天主却召唤了仍关在狱中的\Person{塞昆杜卢斯}，使他更早离世，不无恩宠，以此让野兽有喘息之机。不过，即使他的灵魂没有承认感恩的理由，但他的肉身确实承认了。\par

2. 至于\Person{斐利奇达}（因为主也同样向她施恩），当时她已怀孕八个月（她被捕时已经怀孕），当行刑日临近时，她非常忧愁，唯恐因怀孕被推迟——因为孕妇不允许公开处决——也唯恐她的圣洁无邪的血会与后来作恶的人一起流。此外，她的同伴殉道者也感到痛苦悲伤，生怕留下如此优秀的朋友，如同同伴，独自走在同一希望的道路上。因此，他们同心合意，在行刑前三天向主倾心祈祷。祈祷后，她的产痛立刻来临；当她因八个月分娩的自然困难而阵痛时，一个属于\OriginalTerm{卡塔拉克塔里}{Cataractarii}的仆人对她说：\Quote{你现在这样受苦，当被扔给野兽时，你将如何？你曾拒绝献祭而轻视它们。}她回答说：\Quote{现在我受苦是我自己受苦；但那时另有一位在我内，祂将为我受苦，因为我也将为祂受苦。}于是她生下一个女婴，由某位姐妹收养为自己的女儿。\par

3. 圣神既允许，且藉允许而愿意将这比赛中之事记录成文，尽管我们不配描述如此伟大的荣光，但我们仍遵从那真福佩尔佩图亚的命令，毋宁说是她神圣的托付，再增加一个关于她坚毅与高尚心灵的见证。当护民官因某些诡诈之人的暗示，惧怕她们会被某种魔法咒语从狱中救出，而更加严厉地对待她们时，佩尔佩图亚当面回答他说：\Quote{你为何不至少允许我们得到一点休息，既然我们是要得罪最尊贵的凯撒，且必须在他的诞辰之日作战？若我们届时更加肥壮地被带出来，那岂不正是你的荣耀？} 护民官战栗并脸红，命令更人道地对待她们，以致允许她们的弟兄和其他人进去与她们一同休息；甚至狱卒自己也信任她们了。\par

4. 再者，在前一天，当她们享用那最后一餐——他们称之为自由餐——时，她们并非享用自由的晚餐，而是\OriginalTerm{爱宴}{agape}；她们以同样的坚定向民众说出\Emph{针对他们}的话，宣讲主的审判，为她们殉道的幸福作证，嘲笑那些聚集之人的好奇；同时撒图鲁斯说：\Quote{明天对你们来说还不够，让你们带着愉悦观看你们所憎恨的。今日是朋友，明日是仇敌。但要仔细记下我们的面容，以便在审判之日认出它们。} 于是所有的人都惊奇地离开，因这些事许多人信了。\par

% structure: p0025 source=h2 target=section reason=relative-heading-rank; base=h2

\section{第六章 概述。从监狱他们被欢乐地领进竞技场，尤其是佩尔佩图亚和斐利齐塔。所有人都拒绝穿上不洁的衣裳。他们被鞭打，被扔给野兽。撒图鲁斯两次未受伤。佩尔佩图亚和斐利齐塔被摔倒；他们被召回至萨纳维瓦瑞安门。撒图鲁斯被豹子咬伤，鼓励士兵。他们彼此亲吻，被剑杀死。}

1. 他们胜利的日子来临，他们从监狱进入竞技场，仿佛参加集会，喜乐满面，容光焕发；若有些许畏缩，那也是喜乐，而非恐惧。佩尔佩图亚平静地跟随着，步履姿态如同基督的贵妇，天主所爱的人；她垂下目光，避开大众的注视。再者，斐利齐塔因安全产子而喜乐，以便能与野兽搏斗；从流血和产婆直接到角斗士，并在产后用第二次圣洗来洗。当他们被带到门口，被迫穿上衣服——男子穿萨图尔努斯祭司的服装，女子穿那些献身于刻瑞斯之人的服装——那位高尚的女性坚持到最后一刻。因为她说：\Quote{我们走到这一步是自愿的，正是为此，我们的自由不应受到限制。为此我们已屈服了心神，以免做任何这样的事：我们已与你们协议好了。} 不义承认了正义；护民官让步，允许她们像原本那样被带进来。佩尔佩图亚唱着圣咏，已经践踏了埃及人的头；勒沃卡图斯、撒图尔尼努斯和撒图鲁斯向着围观的人群发出关于此次殉道的威胁。当他们进入希拉里亚努斯的视线时，他们开始以手势和点头对希拉里亚努斯说：\Quote{你审判我们，}他们说，\Quote{但天主将审判你。} 听到这话，民众被激怒，要求他们在经过\OriginalTerm{猎人}{venatores}的行列时受鞭刑。他们确实喜乐，因为他们蒙受了主苦难中的任何一样。\par

2. 但那位说过\BibleQuote{你们求，就必得到}\BibleRef{若 16:24}的主，应允了他们的所求，赐给他们各人所愿的死亡。因为当他们曾彼此谈论他们关于殉道愿望时，撒图尔尼努斯确实宣称他希望被扔给所有野兽；无疑是为了获得更荣耀的冠冕。因此在表演开始时，他和勒沃卡图斯先经历了豹子，此外在台架上还受到熊的骚扰。然而撒图鲁斯最憎恨的莫过于熊；但他想象自己会被豹子一口咬死。因此，当一头野猪被提供时，倒是那个提供野猪的猎人反被那野猪刺伤，并在表演后的第二天死去。只有撒图鲁斯被拖出来；当他被绑在靠近熊的地板上时，熊不肯从洞里出来。于是撒图鲁斯第二次被召回，未受伤害。\par

3. 此外，魔鬼为少年女子准备了一头极其凶猛的母牛，特意为此目的而安排，违反了常规，在野兽方面也与她们的性别较劲。于是，她们被剥去衣服，只披着网，被带了出来。民众看见一位体态纤弱的少女，另一位乳房仍在滴着刚生产后的乳汁，都战栗起来。因此，她们被召回，解开了束缚。佩尔佩图阿首先被领进去。她被抛起，摔在臀部；当她看见自己侧边的长袍被撕破时，便把它拉过来作为遮羞的帕子，心里更挂念她的端庄而非痛苦。然后她再次被叫去，并束起散乱的头发；因为殉道者披头散发受苦是不合适的，免得她在光荣中显出哀悼的样子。于是她站起来；当她看见费利西塔斯被压垮时，便走近她，向她伸出手，扶她起来。两人并肩而立；民众的残暴平息后，她们被召回\Place{桑纳维瓦利亚门}。然后，一位尚是慕道者的人，名叫\OriginalTerm{路斯提库斯}{Rusticus}，接待了佩尔佩图阿，他一直紧跟着她；她仿佛从睡梦中醒来，之前她深深地沉浸在圣神和神魂超拔中，开始环顾四周，对惊讶的众人说：\Quote{我不知道我们何时要被领出去面对那头母牛。}当她听说已经发生的事时，她不肯相信，直到她在自己身体和衣服上看到某些受伤的迹象，并认出了那位慕道者。之后，她让那位慕道者和兄弟走近，对他们说：\Quote{要在信德中站稳，彼此相爱，你们众人，不要因我的受苦而跌倒。}\par

4. 同一位\Person{萨图鲁斯}在另一个入口鼓励士兵\Person{普登斯}说：\Quote{是的，我在这里，正如我曾应许和预言过的，因为直到此刻我没有感受到野兽。现在请全心相信。看哪，我要出去面对那头野兽，我将被豹子一口咬死。}演出一结束，他立刻被扔给豹子；豹子一口咬下去，他全身浸满了鲜血，以致民众在他返回时高声向他呼喊，作为他\OriginalTerm{第二次圣洗}{second baptism}的见证：\Quote{得救又洗净，得救又洗净。}显然，这样光荣地出现在场面中的他确实是得救了。然后他对士兵\Person{普登斯}说：\Quote{保重，并记念我的信德；不要让这些事扰乱你，而要坚固你。}同时，他请求从小指上取下一个小戒指，然后沾满自己的鲜血还给他，留给他一份遗产的信物和他鲜血的纪念。然后他毫无生气地与其他殉道者一起被扔下，要在惯常的地方被处死。当民众呼喊他们到中间去，以便刀剑刺入他们身体时，让他们的眼睛成为杀戮的合伙人时，他们自愿站起来，转移到民众想要的地方；但他们先彼此亲吻，以平安礼来完成他们的殉道。其余的人确实一动不动、沉默地接受了剑刺；\Person{萨图鲁斯}更甚，因为他首先登上梯子，首先交出灵魂，因为他也在等待佩尔佩图阿。但佩尔佩图阿，为了尝受一些痛苦，肋骨之间被刺中，大声呼喊，她自己将年轻角斗士颤抖的右手引向自己的喉咙。也许这样一个妇人除非她自愿，否则不可能被杀，因为不洁的魔鬼惧怕她。\par

最勇敢蒙福的殉道者啊！真正被召选而归于我们的主耶稣基督光荣的人！凡高举、尊敬并朝拜祂的人，确实应当阅读这些榜样，以建立教会，不比古代的榜样逊色，以便新的德能也证明同一位圣神至今仍在运作，全能的天主圣父，以及祂的圣子耶稣基督我们的主，祂享有光荣和无限的能力，直到永远。阿们。\par

\SourceRecord{Translated by R.E. Wallis.；Ante-Nicene Fathers；Vol. 3.；Edited by Alexander Roberts, James Donaldson, and A. Cleveland Coxe.；Buffalo, NY: Christian Literature Publishing Co.,；1885.；<http://www.newadvent.org/fathers/0324.htm>.}
